\documentclass[10pt,oneside]{report}
\usepackage{soc2style}
\begin{document}
\doctitle{SOC2 Training Curriculum}{Twelve-module progression, labs, questions and capstone}{SOC2-08}
\section*{Operating convention and boundaries}
An AICPA criterion identifies an outcome against which controls may be evaluated. The selected point-of-focus headings identify relevant characteristics; they are not discrete certification checklist items. Every owner assignment, technological mechanism and audit test in this suite is a proposed example to be tailored to the system description, commitments, risk profile, period, and service organization responsibilities.

\begin{itemize}
\item Always establish system boundaries, the in-scope trust services categories, service commitments and system requirements first.
\item For a security-only engagement, use the common criteria. For additional trust services categories, include the corresponding additional criteria, subject to properly justified applicability.
\item Document complementary user entity controls and subservice organizations; do not treat vendor assurance reports as replacements for management ownership.
\item Privacy points of focus marked [P] or [C] describe typical processor or controller relevance; assess actual responsibilities rather than treating the annotations as absolute.
\end{itemize}

\chapter{Training architecture}
This curriculum provides a prerequisite-based professional training sequence for security engineers, GRC practitioners, DevSecOps/platform engineers, control owners and internal auditors. The learning outcomes are independent educational recommendations and do not award, imply or replace an AICPA, CPA or SOC 2 credential.

\subsection{Recommended delivery}
Twelve modules, typically one week each, with facilitated seminars, optional technical labs and an integrated capstone. Learners should use the authoritative source criterion table (Document 01), practice cards (Document 02), implementation playbooks (Document 03), and test register (Document 07). Adjust pacing to the audience; the cadence is not a source requirement.

\subsection{Competency ladder}
\begin{itemize}
\item Foundation: accurately locate source criterion, distinguish the TSC from technology-specific advice and explain categories.
\item Practitioner: map service commitments to risks, design control activity and collect sufficient evidence.
\item Engineer: implement controls in IAM, CI/CD, SDLC, cloud, privacy and recovery with reliable observability.
\item Assessor: plan scope-specific testing, challenge population completeness, analyze exceptions and communicate conclusions.
\end{itemize}

\chapter{Prerequisite-ordered modules}
\section{Module 01: Assurance foundations and scope}
Prerequisites: none. Related criteria: CC1.1--CC9.2 (overview); A/PI/C/P categories.

\subsection{Learning objectives}
Distinguish criterion, point of focus, control, audit evidence and Type 1/2; specify scoped system.

\subsection{Hands-on exercise}
Write a one-page system description with three service commitments.

\subsection{Acceptance artifact}
Scope statement and classification quiz.

\subsection{Knowledge check}
\begin{itemize}
\item State the difference between the published control criterion and the technique selected in the lab.
\item Identify which source systems produce evidence and why evidence could be incomplete.
\item Give one realistic failure mode, one compensating control and one question a service auditor should ask.
\end{itemize}

\section{Module 02: Governance, competence and control ownership}
Prerequisites: 01. Related criteria: CC1.1--CC1.5 CC5.3.

\subsection{Learning objectives}
Map oversight, roles, policy and accountability to risks.

\subsection{Hands-on exercise}
Create RACI, board escalation path and control owner roster.

\subsection{Acceptance artifact}
Governance operating model.

\subsection{Knowledge check}
\begin{itemize}
\item State the difference between the published control criterion and the technique selected in the lab.
\item Identify which source systems produce evidence and why evidence could be incomplete.
\item Give one realistic failure mode, one compensating control and one question a service auditor should ask.
\end{itemize}

\section{Module 03: System inventory, objectives and risk}
Prerequisites: 01,02. Related criteria: CC2.1 CC3.1--CC3.4 CC5.1.

\subsection{Learning objectives}
Inventory information assets and create justified risk scenarios.

\subsection{Hands-on exercise}
Map sample SaaS environment and write eight risk responses.

\subsection{Acceptance artifact}
Inventory, DFD and risk register.

\subsection{Knowledge check}
\begin{itemize}
\item State the difference between the published control criterion and the technique selected in the lab.
\item Identify which source systems produce evidence and why evidence could be incomplete.
\item Give one realistic failure mode, one compensating control and one question a service auditor should ask.
\end{itemize}

\section{Module 04: Communication, monitoring and exceptions}
Prerequisites: 02,03. Related criteria: CC2.2 CC2.3 CC4.1 CC4.2.

\subsection{Learning objectives}
Differentiate internal/external communications from control monitoring.

\subsection{Hands-on exercise}
Build defect lifecycle, alert review, communication plan.

\subsection{Acceptance artifact}
Monitoring procedure and findings tracker.

\subsection{Knowledge check}
\begin{itemize}
\item State the difference between the published control criterion and the technique selected in the lab.
\item Identify which source systems produce evidence and why evidence could be incomplete.
\item Give one realistic failure mode, one compensating control and one question a service auditor should ask.
\end{itemize}

\section{Module 05: Identity, access and data movement}
Prerequisites: 03. Related criteria: CC6.1--CC6.7.

\subsection{Learning objectives}
Design joiner/mover/leaver, privilege and data movement controls.

\subsection{Hands-on exercise}
Reconcile users and service accounts; test terminated-user revocation.

\subsection{Acceptance artifact}
IAM control specification and access test.

\subsection{Knowledge check}
\begin{itemize}
\item State the difference between the published control criterion and the technique selected in the lab.
\item Identify which source systems produce evidence and why evidence could be incomplete.
\item Give one realistic failure mode, one compensating control and one question a service auditor should ask.
\end{itemize}

\section{Module 06: Configuration, vulnerability and malicious software}
Prerequisites: 03,05. Related criteria: CC6.8 CC7.1--CC7.3.

\subsection{Learning objectives}
Separate inventory coverage, finding validity, prioritization and security event triage.

\subsection{Hands-on exercise}
Evaluate a mock vulnerability queue and drift alerts.

\subsection{Acceptance artifact}
Evidence-based remediation and monitoring metrics.

\subsection{Knowledge check}
\begin{itemize}
\item State the difference between the published control criterion and the technique selected in the lab.
\item Identify which source systems produce evidence and why evidence could be incomplete.
\item Give one realistic failure mode, one compensating control and one question a service auditor should ask.
\end{itemize}

\section{Module 07: Incident response, recovery and continuity}
Prerequisites: 04,06. Related criteria: CC7.4 CC7.5 CC9.1 A1.1--A1.3.

\subsection{Learning objectives}
Implement containment, recovery, backups and capacity monitoring.

\subsection{Hands-on exercise}
Run tabletop and restore data in safe test environment.

\subsection{Acceptance artifact}
Incident response and recovery exercise report.

\subsection{Knowledge check}
\begin{itemize}
\item State the difference between the published control criterion and the technique selected in the lab.
\item Identify which source systems produce evidence and why evidence could be incomplete.
\item Give one realistic failure mode, one compensating control and one question a service auditor should ask.
\end{itemize}

\section{Module 08: Secure SDLC and third parties}
Prerequisites: 05,06. Related criteria: CC8.1 CC9.2.

\subsection{Learning objectives}
Relate change authorization, secure checks, repository controls and vendor risk.

\subsection{Hands-on exercise}
Create verified PR gate criteria and subservice risk assessment.

\subsection{Acceptance artifact}
Change control evidence pack.

\subsection{Knowledge check}
\begin{itemize}
\item State the difference between the published control criterion and the technique selected in the lab.
\item Identify which source systems produce evidence and why evidence could be incomplete.
\item Give one realistic failure mode, one compensating control and one question a service auditor should ask.
\end{itemize}

\section{Module 09: Confidentiality engineering}
Prerequisites: 03,05. Related criteria: C1.1 C1.2 CC6.7.

\subsection{Learning objectives}
Track classified information through authorized use and disposal.

\subsection{Hands-on exercise}
Create sensitive information lifecycle map and disposal test.

\subsection{Acceptance artifact}
Classification and disposal workpaper.

\subsection{Knowledge check}
\begin{itemize}
\item State the difference between the published control criterion and the technique selected in the lab.
\item Identify which source systems produce evidence and why evidence could be incomplete.
\item Give one realistic failure mode, one compensating control and one question a service auditor should ask.
\end{itemize}

\section{Module 10: Processing integrity and availability}
Prerequisites: 03,07. Related criteria: PI1.1--PI1.5 A1.1--A1.3.

\subsection{Learning objectives}
Define complete, valid, accurate, timely and authorized processing.

\subsection{Hands-on exercise}
Inject invalid inputs; verify output reconciliation and restoration.

\subsection{Acceptance artifact}
Reconciliation and processing-integrity test.

\subsection{Knowledge check}
\begin{itemize}
\item State the difference between the published control criterion and the technique selected in the lab.
\item Identify which source systems produce evidence and why evidence could be incomplete.
\item Give one realistic failure mode, one compensating control and one question a service auditor should ask.
\end{itemize}

\section{Module 11: Privacy data lifecycle and assurance}
Prerequisites: 03,09. Related criteria: P1.1--P8.1.

\subsection{Learning objectives}
Apply controller/processor context, rights requests, notices, consent and disclosure.

\subsection{Hands-on exercise}
Simulate privacy request, processor incident and deletion proof.

\subsection{Acceptance artifact}
Privacy request and incident decision packet.

\subsection{Knowledge check}
\begin{itemize}
\item State the difference between the published control criterion and the technique selected in the lab.
\item Identify which source systems produce evidence and why evidence could be incomplete.
\item Give one realistic failure mode, one compensating control and one question a service auditor should ask.
\end{itemize}

\section{Module 12: SOC 2 readiness and capstone}
Prerequisites: 01--11. Related criteria: All 61 criteria.

\subsection{Learning objectives}
Scope, design, operate, test, and document traceability across a full system.

\subsection{Hands-on exercise}
Conduct simulated readiness assessment and Type 2 evidence sampling.

\subsection{Acceptance artifact}
Complete annotated control matrix and oral defense.

\subsection{Knowledge check}
\begin{itemize}
\item State the difference between the published control criterion and the technique selected in the lab.
\item Identify which source systems produce evidence and why evidence could be incomplete.
\item Give one realistic failure mode, one compensating control and one question a service auditor should ask.
\end{itemize}

\chapter{Scenario-driven assessment bank}
\subsection*{Question 1}
A service provider includes only security in its intended SOC 2 scope. Are A, PI, C and P series automatically included?

\subsection*{Question 2}
An engineer passes CodeQL once, six months before the Type 2 period. Does this prove the check operated during the period?

\subsection*{Question 3}
A developer deletes an exposed password from the default branch. Is rotation optional?

\subsection*{Question 4}
A backup system sends nightly success emails. Can the team claim recovery effectiveness?

\subsection*{Question 5}
An access reviewer signs off on 200 listed users but never compares to the HR termination list. What might be missing?

\subsection*{Question 6}
Are [P] and [C] annotations rigid statutory allocations?

\subsection*{Question 7}
An outsourced data center is responsible for badges. Who must document the relationship?

\subsection*{Question 8}
A monitoring rule creates 500 alerts, none assigned to analysts. Is logging alone conclusive evidence?

\subsection*{Question 9}
A privacy deletion request was honored in production but not in downstream analytics. What failed?

\subsection*{Question 10}
Can a policy document alone prove operating effectiveness?

\subsection*{Question 11}
Is a source-code vulnerability identical to confirmed exploitation?

\subsection*{Question 12}
If a privacy point of focus does not apply to a service processor, may management discard the whole privacy criterion?

\subsection*{Question 13}
Why is output authorization part of processing integrity?

\subsection*{Question 14}
When is accepting a risk distinct from resolving a deficiency?

\subsection*{Question 15}
What is a common reason supplier SOC 2 reports fail to support the entity's claim?

\subsection*{Question 16}
How should an automated control failure during a reporting period be handled?

\subsection*{Question 17}
What does the AICPA say about evaluating every individual point of focus?

\subsection*{Question 18}
Who expresses the independent attestation opinion?

\chapter{Instructor answer key}
\paragraph{Question 1.} No. Security uses common criteria; additional specific series are associated with additional selected categories.

\paragraph{Question 2.} No. Link the period population of relevant changes to successful applicable check results and exceptions.

\paragraph{Question 3.} Evaluate whether the credential was live or exposed; if valid and sensitive, revocation/rotation and misuse evaluation may be necessary. History deletion alone does not prove risk eliminated.

\paragraph{Question 4.} Not on that basis alone. Review restore testing, recovery objectives, integrity and dependencies.

\paragraph{Question 5.} The population may exclude unlisted accounts or terminated users; completeness and deprovisioning must be evaluated.

\paragraph{Question 6.} No. They indicate typical relevance to processors or controllers; actual responsibilities and circumstances govern the analysis.

\paragraph{Question 7.} The service organization still needs to understand and represent relevant subservice responsibilities and how they are addressed in its engagement.

\paragraph{Question 8.} No. Effective monitoring involves analysis, evaluation, routing and response based on the control design.

\paragraph{Question 9.} End-to-end disposal process may be incomplete. Identify all affected data stores and required processor actions.

\paragraph{Question 10.} No. An approved policy may support design; occurrence records and period coverage support operation.

\paragraph{Question 11.} No. A static finding may indicate an issue or potential path; assess runtime reachability, mitigations, context and risk.

\paragraph{Question 12.} Not automatically. Evaluate the criterion, selected category, service commitments and relevant system responsibilities; justify any applicability decision.

\paragraph{Question 13.} Processing integrity includes authorized as well as complete, valid, accurate and timely processing and delivery.

\paragraph{Question 14.} Acceptance documents accountable understanding of residual risk and monitoring; it does not inherently make ineffective control operation effective.

\paragraph{Question 15.} Wrong reporting period, system boundary, carve-out, complementary user entity controls or unresolved report exceptions.

\paragraph{Question 16.} Preserve population and time context, evaluate impact, investigate cause, document exception and remediation, and consult the practitioner on implication.

\paragraph{Question 17.} The source explicitly says assessing whether each point of focus is addressed is not required.

\paragraph{Question 18.} The qualified practitioner/service auditor, not an engineering tool, readiness consultant or generic scoring framework.

\chapter{Capstone specification and rubric}
\subsection{Simulated service context}
Model a fictional business-to-business SaaS platform that processes regulated customer data and proprietary information using source control, container builds, cloud identity, managed databases and third-party observability. Define its actual commitments, selected trust services categories and controller/processor roles explicitly; do not assume all categories are in scope.

\subsection{Required deliverables}
\begin{enumerate}
\item System boundary diagram and trust services category rationale.
\item Risk register with fraud, security, capacity and third-party scenarios.
\item Control register mapped to relevant criteria and owners; documented design and actual simulated implementation evidence.
\item One joined population reconciliation for user lifecycle, changes, vulnerabilities and system recovery.
\item Security incident tabletop, end-to-end recovery test, privacy rights case and documented corrective actions.
\item Evidence manifest, operating test workpaper, exception analysis and executive readout.
\end{enumerate}

\subsection{Illustrative capstone rubric (100 points)}
\begin{itemize}
\item Source fidelity and selected-category completeness -- 20 points.
\item Risk-control reasoning and design consistency -- 20 points.
\item Reproducible population/evidence integrity -- 25 points.
\item Technical exercise effectiveness and exception handling -- 20 points.
\item Clear reporting, accountability and defensible limitations -- 15 points.
\end{itemize}

Suggested completion threshold: 80/100 overall and no critical source-fidelity errors, followed by remediation of identified weak areas. This is a training benchmark only, not an external AICPA assessment requirement.


\chapter*{Source and editorial protocol}
\addcontentsline{toc}{chapter}{Source and editorial protocol}
This independent document is grounded in the user-supplied \emph{AICPA TSP Section 100: 2017 Trust Services Criteria for Security, Availability, Processing Integrity, Confidentiality, and Privacy (With Revised Points of Focus---2022)}, 75-page PDF, especially paragraphs .01--.29, the criterion table on pp.~14--67, and the glossary. Criterion wording is quoted when labeled as such. Selected point-of-focus labels are aids to interpretation and may be incomplete. Engineering recommendations, role assignments, tools, tests, dates, scorecards, and example frequencies are author-developed examples and are not AICPA requirements. The 2022 revisions changed points of focus, not the 2017 criteria (source p.~3). The source states that an assessment of each point of focus is not required (p.~5). This document does not constitute a CPA attestation or legal advice.
\end{document}
